\section{Test and Evaluation Procedures}
\label{sec:ch3-test-evaluation}

The evaluation followed the same order as the implementation: sensor
validation, overload testing, socket-heating characterization,
socket-heating protection validation, series arc-fault testing, and TinyML
model evaluation. Full calibration records are placed in
Appendix~\ref{app:a}; socket-heating characterization and protection data are
placed in Appendices~\ref{app:c} and~\ref{app:d}; arc-fault testing and data
collection logs are placed in Appendices~\ref{app:e} and~\ref{app:f}. The
letters, safety proposal, project schedule, and cost documentation are placed
in Appendices~\ref{app:g},~\ref{app:h}, and~\ref{app:i}.

\subsection{Sensor-Validation Metrics}
\label{sec:ch3-sensor-metrics}

Sensor validation compared prototype readings against reference instruments.
The mean absolute error was computed as
\begin{equation}
  \mathrm{MAE} =
  \frac{1}{n}\sum_{i=1}^{n}|y_i-\hat{y}_i|.
  \label{eq:metric-mae}
\end{equation}
The mean squared error was computed as
\begin{equation}
  \mathrm{MSE} =
  \frac{1}{n}\sum_{i=1}^{n}(y_i-\hat{y}_i)^2.
  \label{eq:metric-mse}
\end{equation}
The root mean squared error was computed as
\begin{equation}
  \mathrm{RMSE} =
  \sqrt{\frac{1}{n}\sum_{i=1}^{n}(y_i-\hat{y}_i)^2}.
  \label{eq:metric-rmse}
\end{equation}
The mean absolute percentage error was computed as
\begin{equation}
  \mathrm{MAPE} =
  \frac{100}{n}\sum_{i=1}^{n}
  \left|\frac{y_i-\hat{y}_i}{y_i}\right|.
  \label{eq:metric-mape}
\end{equation}
In (\ref{eq:metric-mae})--(\ref{eq:metric-mape}), $y_i$ is the reference
measurement, $\hat{y}_i$ is the prototype measurement, and $n$ is the number of
test points. MAE gives the average magnitude of error, MSE and RMSE penalize
larger errors more strongly, and MAPE expresses the error relative to the
reference value [72].

\subsection{Overload Test Procedure}
\label{sec:ch3-overload-test-procedure}

Overload testing used high-power load combinations to determine whether the
prototype warned above the warning threshold and tripped only when the
sustained-overload condition persisted. The block diagram of the overload test
is shown in Figure~\ref{fig:ch3-overload-block}.
\begin{figure}[H]
  \centering
  \resizebox{0.96\textwidth}{!}{%
  \begin{tikzpicture}[node distance=0.34in and 0.36in]
    \node[thesis process,text width=1.35in] (supply) {230V AC supply\\with breaker access};
    \node[thesis process,text width=1.35in,right=of supply] (plug) {TinyML Smart Plug\\under test};
    \node[thesis process,text width=1.35in,right=of plug] (load) {High-power\\load bank};
    \node[thesis process,text width=1.55in,below=of plug] (log) {PWA logs\\current, warning, relay, latency};
    \draw[thesis arrow] (supply) -- (plug);
    \draw[thesis arrow] (plug) -- (load);
    \draw[thesis arrow] (plug) -- (log);
  \end{tikzpicture}%
  }
  \caption{Testing setup block diagram for overload evaluation.}
  \label{fig:ch3-overload-block}
\end{figure}

Figure~\ref{fig:ch3-overload-block} shows the prototype placed between the
supply and load bank so that it could measure current and disconnect the load
through its relay. This matches the current path used in normal operation.

The actual overload test setup is shown in Figure~\ref{fig:ch3-overload-actual}.
\begin{figure}[H]
  \centering
  \includegraphics[width=0.88\textwidth]{ch3/figure_3.31.png}
  \caption{Actual test setup for overload evaluation.}
  \label{fig:ch3-overload-actual}
\end{figure}

Figure~\ref{fig:ch3-overload-actual} documents the physical load arrangement
used to create warning and sustained-overload conditions. During each trial,
mean voltage, mean current, peak current, warning state, relay state, and trip
latency were recorded.

\subsection{Socket-Heating Test Procedure}
\label{sec:ch3-heating-test-procedure}

Socket-heating tests compared normal and intentionally degraded socket
contacts under similar load levels. The objective was to determine whether
degraded contacts produced faster and higher temperature rise, and whether the
external device temperature measurement and compensated estimator could provide useful warning before
dangerous contact heating progressed.

Figure~\ref{fig:ch3-heating-block} shows the socket-heating test block
diagram.
\begin{figure}[H]
  \centering
  \resizebox{0.96\textwidth}{!}{%
  \begin{tikzpicture}[node distance=0.34in and 0.34in]
    \node[thesis process,text width=1.35in] (supply) {230V AC supply\\with breaker access};
    \node[thesis process,text width=1.35in,right=of supply] (plug) {Prototype\\thermistor channel};
    \node[thesis process,text width=1.35in,right=of plug] (socket) {Normal or degraded\\socket contact};
    \node[thesis process,text width=1.35in,right=of socket] (load) {Matched\\load case};
    \node[thesis process,text width=1.55in,below=of socket] (reference) {Infrared thermometer\\terminal/contact reading};
    \node[thesis process,text width=1.55in,below=of plug] (log) {Logs: device temperature measurement, estimate, terminal, relay};
    \draw[thesis arrow] (supply) -- (plug);
    \draw[thesis arrow] (plug) -- (socket);
    \draw[thesis arrow] (socket) -- (load);
    \draw[thesis arrow] (reference) -- (log);
    \draw[thesis arrow] (plug) -- (log);
  \end{tikzpicture}%
  }
  \caption{Testing setup block diagram for socket-heating evaluation.}
  \label{fig:ch3-heating-block}
\end{figure}

Figure~\ref{fig:ch3-heating-block} shows that both the prototype Device
Measurement and a reference socket or terminal temperature measurement were
observed during testing. This made it possible to distinguish device temperature measurement,
estimated socket temperature, and reference terminal temperature in the
results.

The degraded socket contacts are shown in Figure~\ref{fig:ch3-corroded-contacts}.
\begin{figure}[H]
  \centering
  \includegraphics[width=0.80\textwidth]{ch3/figure_3.33.png}
  \caption{Corroded socket contacts used for degraded-contact testing.}
  \label{fig:ch3-corroded-contacts}
\end{figure}

Figure~\ref{fig:ch3-corroded-contacts} documents the degraded-contact
condition used to accelerate contact heating. The degraded socket was compared
against a normal socket under matched load combinations.

The actual socket-heating test setup is shown in
Figure~\ref{fig:ch3-heating-actual}.
\begin{figure}[H]
  \centering
  \includegraphics[width=0.88\textwidth]{ch3/figure_3.34.png}
  \caption{Actual test setup for socket-heating evaluation.}
  \label{fig:ch3-heating-actual}
\end{figure}

Figure~\ref{fig:ch3-heating-actual} shows the bench configuration used to
record time, current, device temperature measurement, estimated socket temperature,
reference terminal temperature, and relay behavior. The characterization data
in Appendix~\ref{app:c} were used to support the thermal model, while the
protection validation data in Appendix~\ref{app:d} were used to evaluate
relay response.

\subsection{Series Arc-Fault Test Procedure}
\label{sec:ch3-arc-test-procedure}

Series arc-fault testing used a controlled arc platform and appliance loads.
The goal was to create repeatable waveform disturbances while recording
current features, model outputs, relay behavior, and event labels. The
assistance letter and safety proposal for this testing activity are included
in Appendix~\ref{app:g}. The arc platform is shown in
Figure~\ref{fig:ch3-arc-platform}.
\begin{figure}[H]
  \centering
  \includegraphics[width=0.92\textwidth]{ch3/figure_3.35.png}
  \caption{Series arc-fault test platform.}
  \figsource{Qi et al., 2025 [48]}
  \label{fig:ch3-arc-platform}
\end{figure}

Figure~\ref{fig:ch3-arc-platform} shows the principle of the series arc test:
the current path is interrupted by an adjustable contact gap, allowing a
controlled arc disturbance to appear in series with the load. The adapted
source describes a 220V/50Hz platform using a fixed electrode and a moving
electrode controlled by a stepper motor [48]. The prototype testing used the
same series-interruption principle as a guide, while the actual bench setup
used in this study is documented in Figure~\ref{fig:ch3-arc-actual}.

The block diagram of the series arc-fault test is shown in
Figure~\ref{fig:ch3-arc-block}.
\begin{figure}[H]
  \centering
  \resizebox{0.96\textwidth}{!}{%
  \begin{tikzpicture}[node distance=0.34in and 0.34in]
    \node[thesis process,text width=1.25in] (supply) {230V AC supply\\with breaker access};
    \node[thesis process,text width=1.28in,right=of supply] (arc) {Series arc\\generator};
    \node[thesis process,text width=1.28in,right=of arc] (plug) {TinyML Smart Plug\\waveform logging};
    \node[thesis process,text width=1.28in,right=of plug] (load) {Appliance\\load};
    \node[thesis process,text width=1.55in,below=of plug] (label) {PWA/Python logs\\features, labels, relay};
    \draw[thesis arrow] (supply) -- (arc);
    \draw[thesis arrow] (arc) -- (plug);
    \draw[thesis arrow] (plug) -- (load);
    \draw[thesis arrow] (plug) -- (label);
  \end{tikzpicture}%
  }
  \caption{Testing setup block diagram for series arc-fault evaluation.}
  \label{fig:ch3-arc-block}
\end{figure}

Figure~\ref{fig:ch3-arc-block} shows the smart plug measuring the load current
while the arc generator is placed in series with the appliance. This setup is
consistent with the feature definitions in Section~\ref{sec:ch3-waveform-features},
which depend on current waveform disturbance.

The actual series arc-fault test setup is shown in
Figure~\ref{fig:ch3-arc-actual}.
\begin{figure}[H]
  \centering
  \includegraphics[width=0.90\textwidth]{ch3/figure_3.37.png}
  \caption{Actual test setup for series arc-fault evaluation.}
  \label{fig:ch3-arc-actual}
\end{figure}

Figure~\ref{fig:ch3-arc-actual} documents the laboratory arrangement used to
collect arc and normal trials. Complete arc-fault test records are referenced
in Appendices~\ref{app:e} and~\ref{app:f}.

\subsection{Classification Metrics}
\label{sec:ch3-classification-metrics}

For arc-fault evaluation, the positive class was the arc-fault condition and
the negative class was normal operation. Accuracy, precision, recall, and
F1-score were computed from true positives (TP), false positives (FP), true
negatives (TN), and false negatives (FN):
\begin{equation}
  \mathrm{Accuracy} =
  \frac{\mathrm{TP}+\mathrm{TN}}
       {\mathrm{TP}+\mathrm{TN}+\mathrm{FP}+\mathrm{FN}},
  \label{eq:metric-accuracy}
\end{equation}
\begin{equation}
  \mathrm{Precision} =
  \frac{\mathrm{TP}}{\mathrm{TP}+\mathrm{FP}},
  \label{eq:metric-precision}
\end{equation}
\begin{equation}
  \mathrm{Recall} =
  \frac{\mathrm{TP}}{\mathrm{TP}+\mathrm{FN}},
  \label{eq:metric-recall}
\end{equation}
and
\begin{equation}
  F_1 =
  2\frac{\mathrm{Precision}\times\mathrm{Recall}}
        {\mathrm{Precision}+\mathrm{Recall}}.
  \label{eq:metric-f1}
\end{equation}
In (\ref{eq:metric-accuracy})--(\ref{eq:metric-f1}), precision measures how
many trip decisions were actually arc faults, recall measures how many arc
faults were detected, and F1-score summarizes the balance between false trips
and missed faults. These metrics are commonly used for supervised
classification and confusion-matrix evaluation [74].
